\documentclass[letterpaper]{article}
\usepackage[utf8]{inputenc}
\usepackage[spanish]{babel}
\usepackage[a4paper]{geometry}
\geometry{top=2cm, bottom=2cm, left=2cm, right=2cm}


\usepackage{pdfpages}
\usepackage{permute}
\renewcommand*\pmtseparator{\ }
\usepackage{amsmath}
\usepackage{bbm}
\usepackage{amssymb}
\usepackage{dashrule}
\usepackage{enumitem}
\usepackage{fancyhdr}
\usepackage{bm}

\pagestyle{fancy}
\fancyhf{} % limpia encabezado y pie


\usepackage[amsmath, thmmarks]{ntheorem}
\theoremseparator{.}
\theoremprework{\hrule}
\newtheorem{eje}{Ejercicio}
\newtheorem{lema}{Lema}

\newtheorem*{sol}{\textbf{Solución}}
\theorempostwork{\vspace{\baselineskip}}
\newtheorem*{dem}{\textbf{Demostración}}

\newcommand{\ep}{\varepsilon}
\newcommand{\fle}{\Rightarrow}
\newcommand{\ssi}{\Leftrightarrow}
\newcommand{\si}{\Leftrightarrow}
\newcommand{\nrm}[1]{\left\lVert #1 \right\rVert}
\newcommand{\ind}{\mathbbm{1}}

\allowdisplaybreaks

\begin{document}
    \begin{flushleft}
        % Logotipo alineado a la izquierda
        \includegraphics[width=0.12\textwidth]{EscudoCimat.png}
    \end{flushleft}
    
    \vspace{-3.2cm} % Ajusta este valor si necesitas que el logo quede más cerca del texto principal
    {\vspace{0.9cm}\hspace{2.2cm}{\LARGE{Introducci\'on a Ciencia de Datos} }\hspace{3cm}{\Large{Tarea 1} }\\
      
      \hspace{2.2cm}
      {\Large
      Fundamentos y Preparaci\'on de los Datos
      }\\

      \vspace{0.5cm}
      \hspace{2.2cm}
      {\large
      Rub\'i Lara \hspace{1cm}
      Eric Palomino \hspace{1cm}
      Eric
      }
    }

% Encabezado: línea, nombre y título
\fancyhead[R]{ -- }
\fancyhead[L]{Fundamentos y Preparaci\'on de los Datos}

% Número de página en el pie, centrado
\fancyfoot[C]{\thepage}


\thispagestyle{empty}


    \begin{enumerate}
      \item \textbf{Distribuci\'on de los residuos estandarizados}
      Bajo el modelo lineal cl\'asico con errores normales, demuestre que los residuos estandarizados
      \begin{align*}
        r_i = \frac{e_i}{\hat{\sigma}\sqrt{1 - h_{ii}}}
      \end{align*}
      tienen, aproximadamente, distribuci\'on $t$ de Student con $n-p-1$ grados de librtad. Explique c\'omo esta propiedad
      justifica su uso en la detecci\'on de outliers.
      \item 
    \end{enumerate}
\end{document}
